\documentclass[12pt]{article}
\usepackage[paper=a4paper,margin=18mm]{geometry}
\usepackage{amsmath,amssymb}
\usepackage{ctex}
\usepackage{xcolor}
\usepackage{enumitem}
\usepackage{tcolorbox}
\usepackage{unicode-math}
\setmainfont{Times New Roman}
\setmathfont{Cambria Math}
\pagestyle{empty}
\definecolor{bg}{HTML}{202124}
\definecolor{fg}{HTML}{F2F4F8}
\definecolor{muted}{HTML}{C9CDD2}
\definecolor{accent}{HTML}{71C7EC}
\pagecolor{bg}
\color{fg}
\setlength{\parindent}{0pt}
\setlength{\parskip}{6pt}
\setlist[itemize]{leftmargin=*,itemsep=5pt,topsep=2pt}
\tcbset{
  colback=bg,
  colframe=accent,
  boxrule=0.5pt,
  arc=2pt,
  left=8pt,
  right=8pt,
  top=6pt,
  bottom=6pt
}
\newcommand{\sectionbox}[1]{%
  \vspace{4pt}
  {\Large\bfseries\color{accent}#1}\par
  \vspace{2pt}
}
\newcommand{\note}[1]{{\color{muted}\small #1}}
\newcommand{\pp}{\,/ /\,}
\begin{document}

{\huge\bfseries 05 双极结型三极管及放大电路基础：公式总结}

\note{根据 ppt/05 PDF 中出现的主要公式整理。公式按课本样式排版，分式均为上下结构。}

\sectionbox{1. BJT 放大状态下的电流关系}

\begin{tcolorbox}
\[
I_E=I_B+I_C
\]
\[
I_C=I_{NC}+I_{CBO}
\]
\[
\alpha=\frac{I_{NC}}{I_E},\qquad
I_C\simeq \alpha I_E
\]
\[
\beta=\frac{\alpha}{1-\alpha},\qquad
\alpha=\frac{\beta}{1+\beta}
\]
\[
I_C=\frac{\alpha}{1-\alpha}I_B+\frac{1}{1-\alpha}I_{CBO}
=\beta I_B+(1+\beta)I_{CBO}
\]
\[
I_{CEO}=(1+\beta)I_{CBO}
\]
\[
I_C\simeq \beta I_B,\qquad
I_B\simeq (1-\alpha)I_E,\qquad
I_E\simeq I_C
\]
\end{tcolorbox}

\sectionbox{2. 放大作用与电压增益示例}

\begin{tcolorbox}
\[
\Delta i_C=\alpha\Delta i_E
\]
\[
\Delta v_o=-\Delta i_C R_L
\]
\[
A_v=\frac{\Delta v_o}{\Delta v_i}
\]
\[
A_v=\frac{0.98\,\mathrm{V}}{20\,\mathrm{mV}}\approx 49
\]
\end{tcolorbox}

\sectionbox{3. 固定偏置共射放大电路的静态工作点}

\begin{tcolorbox}
\[
I_{BQ}=\frac{V_{BB}-V_{BEQ}}{R_b}
\]
\[
I_{CQ}=\beta I_{BQ}+I_{CEO}\simeq \beta I_{BQ}
\]
\[
V_{CEQ}=V_{CC}-I_{CQ}R_C
\]
\[
V_{CE}=V_{CC}-i_C R_C
\]
\[
v_{BE}=V_{BB}+v_s-i_B R_b
\]
\end{tcolorbox}

\sectionbox{4. 饱和、截止与直流负载线常用判据}

\begin{tcolorbox}
\[
I_{CS}\simeq \frac{V_{CC}}{R_C}
\]
\[
\beta I_B<I_{CS}\quad \Rightarrow\quad \text{放大区}
\]
\[
\beta I_B\ge I_{CS}\quad \Rightarrow\quad \text{饱和区}
\]
\[
I_B\simeq 0\quad \Rightarrow\quad \text{截止区}
\]
\[
u_o=V_{CC}-I_C R_C
\]
\end{tcolorbox}

\sectionbox{5. h 参数小信号模型}

\begin{tcolorbox}
\[
v_{be}=h_{ie}i_b+h_{re}v_{ce}
\]
\[
i_c=h_{fe}i_b+h_{oe}v_{ce}
\]
\[
h_{ie}=\left.\frac{\partial v_{BE}}{\partial i_B}\right|_{v_{CE}},
\qquad
h_{re}=\left.\frac{\partial v_{BE}}{\partial v_{CE}}\right|_{i_B}
\]
\[
h_{fe}=\left.\frac{\partial i_C}{\partial i_B}\right|_{v_{CE}},
\qquad
h_{oe}=\left.\frac{\partial i_C}{\partial v_{CE}}\right|_{i_B}
\]
\[
v_{be}\simeq h_{ie}i_b,\qquad i_c\simeq \beta i_b
\]
\end{tcolorbox}

\sectionbox{6. BJT 输入电阻 $r_{be}$}

\begin{tcolorbox}
\[
r_{be}=r_{bb'}+(1+\beta)r_{b'e}
\]
\[
r_{b'e}=\frac{V_T}{I_{EQ}}
\]
\[
V_T\simeq 26\,\mathrm{mV}\quad (T=300\,\mathrm{K})
\]
\[
r_{be}\approx 200\Omega +(1+\beta)\frac{26\,\mathrm{mV}}{I_{EQ}}
\]
\[
r_{be}\approx 200\Omega +(1+\beta)\frac{26\,\mathrm{mV}}{I_{EQ}(\mathrm{mA})}
\quad\text{按课件常用写法}
\]
\end{tcolorbox}

\sectionbox{7. 基本共射放大电路的小信号指标}

\begin{tcolorbox}
\[
v_i=i_b(R_b+r_{be})
\]
\[
i_c=\beta i_b
\]
\[
v_o=-i_c(R_C\pp R_L)
\]
\[
A_v=\frac{v_o}{v_i}
=-\frac{\beta(R_C\pp R_L)}{R_b+r_{be}}
\]
\[
\text{若 } R_b\gg r_{be},\qquad
A_v\approx -\frac{\beta(R_C\pp R_L)}{r_{be}}
\]
\[
R_i=R_b\pp r_{be}\approx r_{be}
\]
\[
R_o\approx R_C
\]
\end{tcolorbox}

\sectionbox{8. 考虑信号源内阻时的源电压增益}

\begin{tcolorbox}
\[
A_{vs}=\frac{v_o}{v_s}
=\frac{v_o}{v_i}\cdot\frac{v_i}{v_s}
=A_v\frac{R_i}{R_s+R_i}
\]
\[
A_{vs}=A_v\frac{R_i}{R_s+R_i}
\]
\end{tcolorbox}

\sectionbox{9. 分压式射极偏置电路的静态工作点}

\begin{tcolorbox}
\[
V_{BQ}\simeq \frac{R_{b2}}{R_{b1}+R_{b2}}V_{CC}
\]
\[
I_{EQ}\simeq \frac{V_{BQ}-V_{BEQ}}{R_E}
\]
\[
I_{CQ}\simeq I_{EQ}
\]
\[
I_{BQ}=\frac{I_{CQ}}{\beta}
\]
\[
V_{CEQ}=V_{CC}-I_{CQ}R_C-I_{EQ}R_E
\simeq V_{CC}-I_{CQ}(R_C+R_E)
\]
\end{tcolorbox}

\sectionbox{10. 分压式射极偏置电路的小信号指标}

\begin{tcolorbox}
\[
R_b=R_{b1}\pp R_{b2}
\]
\[
v_i=i_b\left[r_{be}+(1+\beta)R_E\right]
\]
\[
v_o=-\beta i_b(R_C\pp R_L)
\]
\[
A_v=-\frac{\beta(R_C\pp R_L)}
{r_{be}+(1+\beta)R_E}
\]
\[
R_i=R_b\pp\left[r_{be}+(1+\beta)R_E\right]
\]
\[
R_o\approx R_C
\]
\[
\text{若 } R_E \text{ 被旁路电容交流短路，则 }
A_v\approx-\frac{\beta(R_C\pp R_L)}{r_{be}}
\]
\end{tcolorbox}

\sectionbox{11. 阻容耦合双电源射极偏置电路}

\begin{tcolorbox}
\[
I_E=\frac{V_{EE}-V_{BE}}{R_E}
\]
\[
I_C\simeq I_E
\]
\[
I_B=\frac{I_C}{\beta}
\]
\[
V_{CE}=V_{CC}+V_{EE}-I_C(R_C+R_E)
\]
\[
A_v=-\frac{\beta(R_C\pp R_L)}
{r_{be}+(1+\beta)R_E}
\]
\[
R_i=R_b\pp\left[r_{be}+(1+\beta)R_E\right],
\qquad
R_o\approx R_C
\]
\end{tcolorbox}

\sectionbox{12. 共集电极放大电路（射极输出器）}

\begin{tcolorbox}
\[
V_{CC}=I_{BQ}R_b+V_{BEQ}+I_{EQ}R_E
\]
\[
I_{EQ}=(1+\beta)I_{BQ}
\]
\[
I_{BQ}=\frac{V_{CC}-V_{BEQ}}{R_b+(1+\beta)R_E}
\]
\[
I_{CQ}=\beta I_{BQ}
\]
\[
V_{CEQ}=V_{CC}-I_{EQ}R_E
\]
\[
v_i=i_b\left[r_{be}+(1+\beta)R_L'\right]
\]
\[
v_o=(1+\beta)i_bR_L'
\]
\[
A_v=\frac{(1+\beta)R_L'}
{r_{be}+(1+\beta)R_L'}
=\frac{R_L'}{\dfrac{r_{be}}{1+\beta}+R_L'}\approx 1
\]
\[
R_L'=R_E\pp R_L
\]
\[
R_i=R_b\pp\left[r_{be}+(1+\beta)R_L'\right]
\]
\[
R_o=R_E\pp\frac{R_s+r_{be}}{1+\beta}
\]
\end{tcolorbox}

\sectionbox{13. 共基极放大电路}

\begin{tcolorbox}
\[
A_v=\frac{v_o}{v_i}
=\frac{\beta(R_C\pp R_L)}{r_{be}}
\]
\[
R_i=R_E\pp\frac{r_{be}}{1+\beta}
\]
\[
R_o\approx R_C
\]
\note{共基电路电压增益为正，输入电阻小，输出电阻约为集电极电阻。}
\end{tcolorbox}

\sectionbox{14. 三种基本组态的常用对比公式}

\begin{tcolorbox}
\[
\begin{array}{c|c|c|c}
\text{组态} & A_v & R_i & R_o\\
\hline
\text{共射} &
-\dfrac{\beta(R_C\pp R_L)}{r_{be}} &
R_b\pp r_{be} &
R_C\\[10pt]
\text{共集} &
\dfrac{(1+\beta)R_L'}{r_{be}+(1+\beta)R_L'}\approx 1 &
R_b\pp\left[r_{be}+(1+\beta)R_L'\right] &
R_E\pp\dfrac{R_s+r_{be}}{1+\beta}\\[12pt]
\text{共基} &
\dfrac{\beta(R_C\pp R_L)}{r_{be}} &
R_E\pp\dfrac{r_{be}}{1+\beta} &
R_C
\end{array}
\]
\end{tcolorbox}

\sectionbox{15. 组合放大电路}

\begin{tcolorbox}
\[
A_v=A_{v1}A_{v2}\cdots A_{vn}
\]
\[
R_i=R_{i1}
\]
\[
R_o=R_{on}
\]
\[
A_{v1}=-\frac{\beta_1R_{i2}'}{r_{be1}},
\qquad
A_{v2}=-\frac{\beta_2(R_C\pp R_L)}{r_{be2}+R_E'}
\]
\[
A_v=A_{v1}A_{v2}
\]
\end{tcolorbox}

\sectionbox{16. 复合管等效参数}

\begin{tcolorbox}
\[
r_{be}=r_{be1}+(1+\beta_1)r_{be2}
\]
\[
\beta\simeq \beta_1\beta_2
\]
\[
r_{be}=r_{be1}
\quad\text{（互补复合管常用近似）}
\]
\[
A_v=\frac{(1+\beta)R_L'}
{r_{be}+(1+\beta)R_L'}\approx 1
\]
\[
R_i=R_{b1}\pp R_{b2}\pp\left[r_{be}+(1+\beta)R_L'\right]
\]
\[
R_o=R_E\pp\frac{R_s+r_{be}}{1+\beta}
\]
\end{tcolorbox}

\sectionbox{17. 频率响应}

\begin{tcolorbox}
\[
\left|A_v\right|
=\left|A_{vm}\right|
\frac{1}{\sqrt{1+\left(\dfrac{f_L}{f}\right)^2}}
\frac{1}{\sqrt{1+\left(\dfrac{f}{f_H}\right)^2}}
\]
\[
\left|A_v\right|
\approx
\frac{\left|A_{vm}\right|}
{\sqrt{1+\left(\dfrac{f_L}{f}\right)^2}}
\quad (f\ll f_H)
\]
\[
\left|A_v\right|
\approx
\frac{\left|A_{vm}\right|}
{\sqrt{1+\left(\dfrac{f}{f_H}\right)^2}}
\quad (f\gg f_L)
\]
\[
f=f_L\ \text{或}\ f=f_H
\quad\Rightarrow\quad
\left|A_v\right|=\frac{\left|A_{vm}\right|}{\sqrt{2}}
\]
\[
20\lg\frac{\left|A_v\right|}{\left|A_{vm}\right|}
=-3\,\mathrm{dB}
\]
\end{tcolorbox}

\vfill
\note{生成文件：output/pdf/bjt_formula_summary.pdf}

\end{document}
